% 초록은 본문이 끝난 뒤 각 절의 첫 문장을 이어 붙여 쓴다(과제 안내 5/5).
% 수치는 모두 generated/results.tex 매크로에서 온다. 숫자를 손으로 적지 않는다.
\begin{abstract}
\noindent
\textbf{문제의식.} 게임과 앱은 ``곧 무언가 나온다''는 시각 효과를 자주 쓰지만, 그런 효과가 사용자의 반응을 실제로 빠르게 하는지는 대개 확인 없이 쓰인다.
실험실 연구에서는 목표 직전의 중립적 경고 신호가 단순 반응시간을 줄인다고 알려져 있다.
이 연구는 그 결과가 참가자 각자의 휴대폰·PC에서, 감독 없이 하는 짧은 웹 과제에서도 나타나는지 확인했다.
\textbf{방법.} 직접 만든 웹 실험 페이지에서 참가자는 화면 가운데 원이 나타나면 최대한 빨리 눌렀다.
절반의 시행에서는 원이 나오기 500\,ms 전에 화면 테두리가 100\,ms 동안 밝아졌고(경고 신호 있음), 나머지 절반에서는 아무 변화가 없었다.
두 조건은 한 참가자 안에서 시행마다 무작위로 섞였다(1인당 본 시행 40회).
분석 규칙은 데이터 수집 전에 고정해 공개 저장소에 기록했다.
\textbf{결과.} \CollectStart{}부터 \CollectEnd{}까지 \Nrecruited{}명이 참여했고, 사전에 정한 제외 기준을 적용한 뒤 \Nanalyzed{}명을 분석했다.
참가자별 반응시간 중앙값의 차이(신호 있음 $-$ 없음)는 중앙값 \DiffMedian{}\,ms(부트스트랩 95\% CI [\BootCIlo, \BootCIhi])였고,
사전 등록한 단측 Wilcoxon 부호순위 검정은 $p$ \Wp{}였다.
\Nnonzero{}명 중 \Kfaster{}명이 신호가 있을 때 더 빨랐다.
\textbf{결론.} 사전에 정한 판정 규칙에 따라 가설은 \ResVerdict.
\ifnum\ResSupported=1
화면 테두리 점등처럼 UI 효과에 가까운 신호도, 통제되지 않은 개인 기기 환경에서 같은 사람의 반응을 빠르게 했다.
다만 이 설계는 신호가 주는 각성과 ``언제 나올지''에 대한 시간 정보를 분리하지 못한다.
\else
이 결과는 ``효과가 없다''가 아니라 이 표본과 환경에서 효과를 확인하지 못했다는 뜻이며, 이 표본이 잡아낼 수 있는 최소 효과크기는 $d_z \approx$ \MDEdz{}였다.
\fi

\medskip\noindent
\textbf{주제어:} 경고 신호, 각성, 단순 반응시간, 웹 실험, 참가자 내 설계, 사전 등록
\end{abstract}
